% QMB_n03_verschraenkung.tex — Kapitel 3: Verschränkung und scheinbare Instantanität
% Inhaltliche Grundlage: Dok. 131, 230, 334, 364; eigene Darstellung
\chapter{Verschränkung ohne Fernwirkung und scheinbare Instantanität}
\label{ch:verschraenkung}

\section{Was Verschränkung in der FFGFT ist}

Zwei Qubits heißen verschränkt, wenn ihr gemeinsamer Zustand nicht als
Produkt $|\psi_1\rangle\otimes|\psi_2\rangle$ geschrieben werden kann.
Der maximal verschränkte Bell-Zustand
\begin{equation}
  |\Phi^+\rangle = \tfrac{1}{\sqrt2}(|00\rangle+|11\rangle)
\end{equation}
hat in der FFGFT eine geometrische Charakterisierung: Die beiden Moden
teilen einen gemeinsamen Träger $T^4/\mathbb{Z}_3$. Ihre
$(z_i,r_i,\theta_i)$-Koordinaten sind nicht unabhängig --- sie sind
Knoten derselben topologischen Konfiguration, die nicht auf Produktform
reduziert werden kann.\status{B}

\textbf{Verschränkung ist in der FFGFT eine topologische Verbindung,
keine nichtlokale Korrelation.} Die Teilchen hatten nie zwei getrennte
Träger; sie teilen von Anfang an denselben. Was die QM als Fernwirkung
beschreibt, ist eine Eigenschaft der gemeinsamen Geometrie.

\section{EPR und Lokalität}

Das EPR-Paradoxon (1935): Wenn QM vollständig ist, sind verschränkte
Teilchen nichtlokal verbunden; Einstein hielt das für inakzeptabel und
vermutete verborgene Variablen.\status{E}

Die FFGFT löst das Paradoxon geometrisch. Lokalität gilt auf dem Torus
in dem Sinn, dass kein Signal schneller als Licht übertragen wird. Die
Korrelation ist eine Eigenschaft der gemeinsamen Geometrie, keine
Signalübertragung. Was aufgegeben wird, ist nur die Vorstellung, dass
jedes Teilchen einen \emph{unabhängigen} Träger hätte --- und die ist
nicht beobachtbar.\status{B}

\section{Scheinbare Instantanität: kein physikalisches Problem}

Der Kollaps der Wellenfunktion scheint instantan zu geschehen: Man misst
Qubit A und B weiß sofort, in welchem Zustand es ist, egal wie weit
entfernt. Das scheint der Relativitätstheorie zu widersprechen.

Bevor man die FFGFT-Auflösung erklärt, lohnt ein Schritt zurück. In
der Standardmathematik gelten Rechenoperationen als instantan: Wenn
$2 + 2 = 4$, dann gilt das gleichzeitig für alle Stellen des Ausdrucks
--- niemand verknüpft damit eine physikalische Instantanität, die
kausal problematisch wäre. Die Gleichheit ist eine strukturelle
Aussage über das Objekt, kein Ereignis im Raum. Genauso verhält es
sich mit der Zustandsbeschreibung in der QM: Die Wellenfunktion
$|\Phi^+\rangle = \frac{1}{\sqrt2}(|00\rangle+|11\rangle)$ ist eine
mathematische Beschreibung eines Zustands --- kein Prozess, der sich
irgendwo ausbreitet. Die Instantanität der Beschreibungsanpassung nach
der Messung ist mathematischer, nicht physikalischer Natur.

In der FFGFT gilt dasselbe Argument auf der Ebene des Trägers: Die
Relation $\tilde T\cdot m = 1$ ist eine \textbf{strukturelle
Eigenschaft} jeder Mode auf $T^4/\mathbb{Z}_3$ --- analog dazu, dass
eine Gleichung für alle ihre Variablen gleichzeitig gilt. Sie ist eine
\textbf{lokale Zwangsbedingung} am selben Raumpunkt, kein
Signal zwischen entfernten Punkten.

Die Auflösung kommt aus der Trennung zweier Ebenen: Eine Änderung bei $\mathbf{x}_1$ breitet sich dann über
die dynamischen Feldgleichungen aus:\status{B}
\begin{equation}
  \text{Änderung bei }\mathbf{x}_1
  \to \text{Ausbreitung mit }v\le c
  \to \text{Wirkung bei }\mathbf{x}_2.
\end{equation}
Die Zeitverzögerung ist $\Delta t = |\mathbf{x}_2-\mathbf{x}_1|/c$.

Lokale Anpassung und globale Ausbreitung sind zwei verschiedene Prozesse.
Die lokale Anpassung ($T\cdot m = 1$ am selben Punkt) erfolgt auf der
Planck-Zeitskala $t_P\sim10^{-43}$\,s --- praktisch instantan für jede
messbare Zeitskala. Die Ausbreitung zu entfernten Punkten ist durch die
Lichtgeschwindigkeit begrenzt. Diese Hierarchie erklärt, warum viele
Quantenprozesse scheinbar instantan wirken.

\section{Die Feldgleichung und ihre Kausalstruktur}

Die FFGFT-Felddynamik folgt einer Wellengleichung:
\begin{equation}
  \frac{\partial^2\tilde T}{\partial t^2} = c^2\nabla^2\tilde T + Q(\tilde T,m,\rho),
\end{equation}
mit der Quellfunktion $Q = -4\pi G\rho\cdot\tilde T$. Das ist eine
hyperbolische Differentialgleichung, charakteristisch für
Wellenausbreitung mit endlicher Geschwindigkeit.

Die Kausalstruktur ist durch die Greensche Funktion manifestiert:
\begin{equation}
  G(\mathbf{x},\mathbf{x}',t-t') = \theta(t-t')\cdot
  \frac{\delta(|\mathbf{x}-\mathbf{x}'|-c(t-t'))}{4\pi|\mathbf{x}-\mathbf{x}'|}.
\end{equation}
Die Heaviside-Funktion $\theta(t-t')$ garantiert Kausalität: Keine
Wirkung vor ihrer Ursache. Die Delta-Funktion kodiert Ausbreitung mit
Lichtgeschwindigkeit. Diese Struktur ist vollständig mit der speziellen
Relativitätstheorie kompatibel.\status{B}

Ein konkretes Beispiel: Ändert sich $m(\mathbf{x}_0)$ bei $t = 0$, dann
passt sich $\tilde T(\mathbf{x}_0)$ sofort an (lokale Zwangsbedingung);
gleichzeitig entsteht $\nabla^2\tilde T\ne0$, das eine Feldstörung
erzeugt; diese breitet sich als Welle mit $v = c$ aus und erreicht
$\mathbf{x}_1$ zur Zeit $t = r/c$.

\section{Die Hierarchie der Zeitskalen}

Die scheinbare Instantanität resultiert aus der extremen Trennung
verschiedener Zeitskalen:

\begin{center}
\small\renewcommand{\arraystretch}{1.3}
\begin{tabular}{p{3.5cm}p{2.5cm}p{3.2cm}}
\toprule
\textbf{Prozess} & \textbf{Zeitskala} & \textbf{Charakter} \\
\midrule
Lokale T-E-Anpassung & $t_P\sim10^{-43}$\,s & unmessbar schnell \\
Wellenfunktionsevolution & $\sim10^{-15}$\,s & Quantenbereich \\
Kausale Feldausbreitung & $t = r/c$ & relativistisch \\
Makroskopische Messung & $\gg10^{-9}$\,s & klassisch \\
\bottomrule
\end{tabular}
\end{center}

Prozesse auf der Planck-Skala sind so schnell, dass sie mit keiner
vorstellbaren Technologie zeitlich aufgelöst werden könnten. Was als
instantan erscheint, läuft auf dieser Skala --- und ist damit kausal
vollständig verträglich.

\section{Die vollständige Dualität: Zeit, Masse, Energie und Länge}

Die FFGFT beschreibt ein umfassendes System von Dualitäten, in dem alle
fundamentalen Größen durch $\tilde T\cdot m = 1$ miteinander verbunden sind:
\begin{align}
  \tilde T\cdot E &= 1\qquad\text{(Zeit-Energie-Dualität)} \\
  \tilde T &= 1/m\qquad\text{(Zeit-Masse-Beziehung)} \\
  E &= mc^2\qquad\text{(Masse-Energie-Äquivalenz)} \\
  \ell &= \hbar/(mc)\qquad\text{(Compton-Wellenlänge als Länge-Energie-Brücke)}
\end{align}
An der Planck-Skala konvergieren alle diese Dualitäten; dort sind Zeit,
Länge, Masse und Energie nicht mehr unabhängig:
\begin{equation}
  t_P = \ell_P/c = \sqrt{\hbar G/c^5}\approx5{,}4\times10^{-44}\,\text{s}.
\end{equation}
Die Dualitäten sind nicht nur mathematische Formulierungen, sondern
haben physikalische Konsequenz: Sie zeigen, dass die scheinbar
verschiedenen Konzepte von Zeit, Raum, Masse und Energie verschiedene
Aspekte derselben geometrischen Realität sind. Je nach Skala dominiert
ein anderer Aspekt: Bei atomaren Skalen die Zeit-Energie-Dualität, bei
makroskopischen Skalen die Masse, bei kosmologischen die Länge-Zeit-Relation.

\section{Topologische Verbindung: der heutige Stand}

Die frühen FFGFT-Dokumente (Herbst 2025) versuchten, die Bell-Korrelationen
durch eine $\xi$-Dämpfung zu erklären, die den CHSH-Wert leicht unter
die Tsirelson-Schranke drückt. Dieser Ansatz ist durch Dok.~230 (Mai 2026)
überholt. Die heutige Position:

Verschränkte Teilchen teilen einen Träger $T^4/\mathbb{Z}_3$ ohne
Fernwirkung. Die statistischen Vorhersagen sind identisch mit der
Quantenmechanik.\status{B} Es gibt keine Fernwirkung, weil es keinen
getrennten Träger gibt. Die topologische Verbindung ist keine
Wirkung über Distanz, sondern eine geometrische Tatsache des gemeinsamen
Trägers. Eine kleine $\xi$-Abweichung $\Delta\mathrm{CHSH}\approx
\xi/(2\pi)\approx10^{-5}$ pro Messung ist als Arbeitshypothese offen,
aber nicht mehr das Kernargument.\status{S}

\begin{laierspur}
Zwei Punkte auf demselben Torus sind verbunden --- nicht durch ein
Signal zwischen ihnen, sondern weil sie dieselbe Oberfläche teilen. Das
ist der Unterschied: EPR fragte nach einer Verbindung zwischen zwei
getrennten Objekten. Die FFGFT sagt: Die Prämisse ist falsch. Die
Objekte waren nie getrennt.
\end{laierspur}

\quelle{Dok.~131 (Scheinbare Instantanität, 2026), Dok.~230 (Verschränkung als topologische Verbindung, Mai 2026), Dok.~334, Dok.~364}
